\chapter{UT Computation Center}

The evolution of UT Austin into a home for computational modeling and simulation is centered in the two decades from 1958 through the end of the seventies. The Computation Center drove this growth by bringing industry expertise onto campus and jump starting a tradition of pioneering research in numerical computing and systems software. When David Young arrived in the fall of 1958, he brought a second thread into the story of early computing at UT Austin. There was an earlier thread already underway, from 1955 at the latest, centered on Al Matsen and his work in the chemistry department with an IBM 650, and both threads are of special interest because of their links with the subsurface and the space race. There was Exxon Houston’s 1958 donation of its IBM CPC to Matsen and the chemistry department in Welch Hall, representing the pragmatism of the oil industry and its ties with Matsen. And there was the 1960 purchase of a CDC 1604 for the math department in Benedict Hall on the South Mall, representing the space race systems thinking of the aerospace industry and its ties with David Young. 

One of the interesting aspects of this very early period in the fifties is the clear differentiation between IBM and CDC hardware. There’s no question that at the time Exxon and the oil industry were using IBM hardware and that this heavily influenced the chemistry department at UT, which also ended up having two IBM machines, both the 650 and the CPC. Meanwhile, David Young was associated with TRW, which was a UNIVAC shop. CDC was a spinoff from UNIVAC, and when Young arrived at UT, CDC hardware followed soon after. There was clearly an important contrast between the oil and aerospace industries at work here, and between the nature of IBM and CDC and their customers.

Both Matsen and Young are variously described as founders and first directors of the UT Computation Center from 1958 up through roughly 1970, and it seems likely that they collaborated in the Computation Center’s early days and blended together the influences of their respective industries and technical fields. A curious fact is that when the CDC 6600 arrived in its underground home in 1966, it was located roughly midway between Welch Hall to the north and Benedict Hall to the south. The 1966 Computation Center sub-terrace building and multi-million dollar 6600 were a landmark for the end of the early days, symbolizing the onset of computing maturity, and how computers and software had grown larger than particular industries and departments.

\section{IBM CPC At Exxon Precursor Humble Oil And UT}

The first concrete event was the chemistry department’s 1955 acquisition of an IBM 650 Magnetic Drum Data Processing Machine using specific research grant funds secured by Al Matsen. The 650 was in Welch, the chemistry building, and was the first real computer on campus. It was a very early mass-produced computer, and Matsen and colleagues used it to compute the Quantum Chemistry Integrals and Tables, which provided the computational foundation for molecular orbital calculations across the field. Although the machine was purchased for his own work, Matsen established a precedent of shared usage at the university by allowing faculty members and researchers from other academic disciplines to utilize the computer. In one legendary exchange, the university president complained to Matsen that the "computer center" did not have long enough open hours and help was not always available. Matsen informed the president that UT actually had no official computer center and was merely using his grant-funded machine, but emphasized that the university desperately needed to build a centralized facility.

A second significant event occurred in 1958 when Exxon precursor Humble Oil in Houston donated its IBM Card-Programmed Electronic Calculator to the university. This was the machine that Rachford, Peaceman, and Douglas had created ADI on four years earlier, in 1954. Matsen was a consultant for Exxon Houston and New Jersey for over thirty-five years. In his Reminiscences he relates a story “Amusingly, I had been lecturing at an unnamed university on the unitary group formulation of the many-body theory. I apparently went way over the listeners' heads since the only question I got was, What possible use could you be to Exxon?” The CPC was a landmark gift and a direct result of Matsen’s extensive ties. To bypass bureaucratic paperwork, Matsen, his graduate students, and other faculty physically carried the heavy machine components into Welch and installed it themselves.

The IBM CPC was not a modern computer, and in fact was a major step backwards from the IBM 650, but it was useful in the Welch Hall of 1958. The CPC will always be legendary as the machine that George Dantzig implemented the Simplex method on at RAND in 1952, and was directly descended from the IBM machines that Richard Feynman worked with at Los Alamos in 1944. It was a hybrid electro-mechanical system. It consisted of an IBM 402 or 417 Accounting Machine controller connected to an IBM 604 Electronic Calculating Punch arithmetic unit and an electromechanical storage unit. It functioned as a decentralized network of specialized units rather than a unified stored-program architecture and was fundamentally incapable of holding both data and instructions. Consequently, a program existed not as a digital state within the machine, but as a physical sequence of punched cards. This required the operator to function as a manual control unit, physically re-entering card decks to execute the iterative loops of a program.

In 1958 David Young moved to UT from TRW and the Los Angeles aerospace environment. He was tasked with founding the Computation Center and serving as its first director. When Young arrived, the Computation Center was "almost non-existent," consisting of Young, his colleague Robert Gregory, and a secretary sharing a single office next to the IBM 650. Within 18 months, Young leveraged his formidable reputation to secure a \$400,000 NSF grant for the CDC 1604, and in 1966, he secured the first \$1,000,000 NSF grant towards the purchase of the CDC 6600 supercomputer. Emphasizing that this was a personal triumph for Young rather than a political favor, Gregory stated: "I could not have done this, you could not have done this, but David Young did it."

I joined David at UT six months after he arrived there in Fall 1958. At that time the Computation Center was almost non-existent. They had acquired an IBM 650 and David, I, and a secretary shared a single office next to the computer room. Within 18 months David, on the basis of his reputation alone, got the first \$400,000 grant from NSF towards the purchase of a CDC 1604, the first transistorized computer. It beat the IBM 7090 into production by one month. Thus, UT went from nothing to a first class Computation Center in one big jump. Then in 1966, after acquiring a building to house the CDC 1604, and after it became saturated with users, David (again on his reputation alone) got the first \$1,000,000 grant from NSF towards the purchase of a CDC 6600. This, again, put UT at the front of the line as far as University Computing Centers go. UT was the first university to have a 6600 and this put them ahead of Berkeley, Stanford, MIT, Harvard and all the rest.

With David Young's TRW background, the center’s institutional goals naturally focused on creating robust numerical frameworks and establishing engineering and scientific computing as an independent discipline. The Computation Center was a catalyst for growth and an environment where academic research converged with government-funded projects. Young’s own career exemplified this duality, as his work on iterative methods for linear systems combined aerospace engineering with applied mathematics. He oriented the Computation Center towards faculty and researchers involved in computer-based research and education, and by 1970 it had evolved into a world-class laboratory with a focus on mathematical programming and a staff with joint appointments in mathematics and computer science that was becoming woven into the university's academic culture. Young recognized that fully exploiting the CDC 6600's hardware required equivalent software innovations. This realization drove him to step down as the Computation Center director in 1970 to found the Center for Numerical Analysis, a name equivalent with the Center for Computational Linear Algebra. His research focused on the software side of the supercomputing equation, developing advanced, highly optimized iterative algorithms such as the Itpack software suite, specifically engineered to solve massive systems of equations on high-performance, parallel architectures like the 6600.

\section{CDC 6600}

The defining characteristic of the Computation Center in the sixties was its aggressive pursuit of the most advanced hardware for floating-point numerical computing. Already from 1960 onward David Young had tied the university to CDC and the numerical-focused Minnesota computer tradition. In the summer of 1966, the university made a landmark investment of six million dollars to acquire a CDC 6600 supercomputer. At the time of its installation, it was the fastest computer in the world, having surpassed the IBM 7030 Stretch by a factor of three. The acquisition of the 6600 elevated UT into an exclusive group, establishing a computational environment that rivaled national facilities like Los Alamos, Livermore, and CERN. For example, CERN installed a 6600 in 1965 and it served as a major computational workhorse for analyzing complex physics data from CERN particle accelerators. The 6600 not only redefined the university's technical capabilities but also fundamentally shaped its intellectual culture, demanding a deep focus on architectural optimization and parallel processing.

A special aspect of the CDC 6600, designed by Seymour Cray and James Thornton and introduced in 1964, lay in its strict division of labor, which was a radical departure from the single-processor mainframes of the era. By offloading all external operations and operating system housekeeping to ten peripheral processors, the central processor was freed to dedicate its 60-bit word length to high-speed arithmetic. The peripheral processors were essentially independent computers based on the earlier 12-bit CDC 160 design. Beyond the distributed external architecture, the 6600 also pioneered internal parallel processing within the central processor itself. The processor was divided into ten distinct, parallel functional units, such as floating-point multiply, divide, add, and shift, allowing the machine to execute multiple instructions simultaneously. Supported by an instruction stack, an early form of instruction cache, and a scoreboard mechanism to track resource availability, the 6600 could issue a new instruction every hundred nanosecond clock cycle, pushing its performance to roughly three million floating-point operations per second, three times faster than the competition. This distributed architecture was a major innovation. The peripheral processors were essentially ten smaller computers that acted as a logistics service. They managed the noisy and slow tasks, such as reading punched cards, spinning magnetic tapes, and updating the console, and freed the main processor to dedicate its entire capacity to high-speed floating-point arithmetic. By delegating input and output operations to ten independent peripheral processors, the 6600 bypassed bottlenecks that had hampered earlier generations of mainframes.

This infrastructure not only provided the raw speed necessary for numerical modeling and simulation but also institutionalized a culture of hardware optimization, where the efficiency of software was measured by its ability to harness the 6600’s unique multi-processor environment. Parallel and multi-processing were encoded into the Computation Center’s genes early on, in an era when simpler single-processor architectures were predominant. This made the environment and technical culture at UT special, and more closely related to the national laboratories such as Los Alamos and Livermore than usual for a university.

The physical reality of the CDC 6600 permanently altered the university's architectural landscape. Because the supercomputer required freon-based cooling systems and large scale power infrastructure, the university constructed a specialized facility, the underground Computation Center situated below the terrace to the east of the Main Building. This unusual home served as the central hub for a growing, decentralized network of remote job entry stations and interactive time-sharing terminals distributed across campus buildings like Pearce Hall and Taylor Hall. While this underground placement preserved the iconic sightlines of the campus, it also addressed the 6600’s formidable infrastructure needs. The CDC 6600 simply could not be supported by the wooden floors and older infrastructure of earlier departmental buildings. The system contained 400,000 individual transistor components and more than a hundred miles of internal wiring. Seymour Cray’s obsession with minimizing signal delays dictated the 6600's unique form factor. Because the central processor was designed to run at a then very fast ten megahertz clock speed, Cray realized that the physical distance between components would cause signaling delays that could limit performance, which led him to pack the 6600's components as tightly as possible to minimize time-of-travel for electrical signals. 

The logic modules, each containing sixty-four silicon transistors packaged in a dense "cordwood" construction, were sandwiched directly between two aluminum cold plates. This density created a heat problem that air cooling could not solve. The 6600 used a freon refrigerant system, the same technology found in sixties air conditioners, but channeled through cold plates directly in contact with the chassis. The refrigerant circulated through the outermost eighteen inches of the machine's four arms, conductively pulling heat from the cold plates and exchanging it with an external chilled water supply. Being underground provided a natural thermal sink, helping the massive refrigeration system manage the heat generated by the transistors. The wiring was just as critical. Cray famously insisted that no wire in the 6600 be longer than eight feet, forcing the machine into its iconic plus-sign or four-wing shape. This ensured the clock signal could reach every component simultaneously. 

This period marked the beginning of a pattern and a special relationship with the Minnesota based Control Data Corporation and its successor, the Cray Research corporation, with the university aligning itself with the same technical lineage that powered the national laboratories. Three generations of CDC machines, the 1604, 6600, and Cyber, preceded the first Cray machine. In the seventies, the university established a similar relationship with DEC, and the Computation Center notably concentrated on CDC and DEC hardware, with the two old IBM systems in Welch from the fifties being the last of their kind for engineering and scientific purposes.

The 6600 was explicitly designed for use at the national labs and in defense and classified projects, and marketed as a supercomputer, a specialized category of machine designed for engineering and scientific tasks. The initial requirements and demand for the CDC 6600 was driven almost exclusively by Los Alamos, Livermore and the NSA. For example, the need at the national labs to simulate shockwaves involving variables changing by twelve orders of magnitude in microseconds made 36-bit floating-point inadequate. On an IBM 7094, simulating these gradients in single precision resulted in numerical noise that obscured the physics. Switching to double precision via software emulation, however, crippled the machine's speed. The 60-bit word was the pragmatic solution to this and other impasses, so that specific requirements led directly to the word size and significantly increased floating-point precision and robustness. One perspective is that 60-bits was the minimum required to do simulations without slowing down for software emulation of double precision, using two words instead of one. Essentially, a 60-bit word with 11-bit exponent and 48-bit fraction allowed what would be double precision on a 36-bit machine to run as single precision on the 6600. With a 48-bit fraction, the machine offered approximately fourteen decimal digits of precision, nearly double the eight digits available on 36-bit machines. 

The inclusion of specialized bit patterns for indefinite and infinity provided a robust error-handling framework that prevented the silent propagation of mathematical anomalies, a precursor to the modern NaN or not a number representation. The 6600 was one of the first machines to emphasize these numerical concepts, and both anticipated and directly influenced features of the IEEE 754 standard two decades later. Unlike modern two’s complement systems, the 6600 used one’s complement arithmetic. This resulted in both a positive zero and a negative zero. While it added complexity for programmers, it allowed for elegant hardware symmetry, since the same comparison logic used for integers could be applied to floating-point values by simply complementing the exponent. This was a classic kind of idiosyncrasy dictated by the desire to optimize numerical performance, as added complexity on one hand allowed for a unified hardware logic that simplified comparisons across different data types on the other hand. While unusual even for its day, this choice became an architectural tradition for Control Data Corporation, with successors continuing to use ones' complement hardware until the late eighties. 

There were more idiosyncrasies to the 6600, especially around non-numerical processing, in other words characters, text, and what could be termed business-style computing, as opposed to scientific-style computing. Notably, IBM hardware would appear in the business school and in the administrative sections of the Main Building. The CDC 6000 series architecture did not have byte addressability and characters were stored as 6-bit codes, with ten characters packed into each 60-bit memory-word. To change a single letter in the middle of a memory word, a programmer had to load the word into a register, use masking and shifting operations to isolate the six bits, swap them, and return the word. A technical divide existed in the CDC community between users of a sixty-four character set and users of a so-called sixty-three character set. The sixty-four character set was the CDC default, and used null-bytes and a full zero word to signal the end of a text string. 

UT adopted the sixty-three character set, which curiously enough was also popular in the Netherlands. In this system, line termination was signaled by two or more null-bytes at the end of a single 60-bit word. This choice reflected the university's independent spirit in software design and its willingness to deviate from vendor standards to achieve what was perceived as better internal efficiency. This sixty-three character culture meant that software often required tweaking to handle the local character set and end-of-line signals. The lack of byte addressability and the 6-bit character limitation forced a generation of programmers to become masters of low-level optimization. This culture of packing and shifting and the decision to maintain a sixty-three character set were deliberate acts of technical independence. They allowed the university to build an ecosystem that was highly optimized for its specific research needs, even if it meant more work when porting software from other institutions. This all played a role in building a supercomputing heritage at UT. The Computation Center's organizational structure, the tripartite division of operations, systems, and analysis, provided the blueprint for what would eventually become the Center for High Performance Computing followed by the Texas Advanced Computing Center from around the late nineties.

\section{The UT2D Dual Dinosaur}

Harnessing the multi-layered parallel architecture of the 6600 required the Computation Center to cultivate a highly specialized software environment. CDC’s early operating systems were often plagued by delays or inefficiencies, and local staff developed the UT2D operating system to manage the specific needs of the university community. This system eventually orchestrated a dual-mainframe complex linking the 6600 with a later 6400, a less expensive model that used a serial central processor instead of the 6600's parallel functional units. The 6600 and 6400 were used together in an integrated way that led to extensive local systems programming. Here is a description of the system from a 1975 research paper.

\begin{quotation}
The Computation Center at The University of Texas at Austin campus provides computing services for both students and researchers. The Computation Center is equipped with a Control Data Corporation (CDC) Model 6600 computer system, which was purchased and installed in 1966, and a CDC model 6400 computer, which was purchased and installed in 1971. The system (UT2D) is one of the most extensive academic computing systems in the United States. UT2D runs 24 hours a day, 7 days a week. High-speed card readers and paper printers are situated at 16 strategic locations on campus, three of which are in the engineering building complex. A timesharing system is provided through the Bell telephone system, allowing 128 users simultaneous connection to the central computer via keyboard terminals, teletypewriters, or similar apparatus. At the keyboard terminal, the user can write programs and run them in.a conversational or interactive mode. The UT2D system includes permanent files which allow the user to reference data, programs, and binary decks without reconstructing a file each time he references it in successive jobs. \cite{derrSystem2000Data}
\end{quotation}

The standard operating system provided by CDC was Scope, a batch-oriented system that many academic institutions found too restrictive for the varied needs of thousands of students and diverse research faculty. The Computation Center systems group focused on modifying and creating software that could maximize the 6600’s throughput. This led to the adoption of Mace, regarded as faster and more configurable than Scope, and it served as the base for many of the university's local modifications throughout the late sixties. The earlier Scope and Mace systems led onward to UT2D, nicknamed the Dual Dinosaur, an operating system fully developed in-house by the Computation Center's staff. UT2D was designed with an innovative job scheduler. If a student submitted a low-cost, ten second compilation job, the system would preemptively pause a larger, high-cost simulation, run the student's code, and then immediately resume the larger job. This was in direct alignment with the university's educational mission. Because student jobs, typically for compiling and running smaller programs, tended to be low-cost, the system automatically gave them the highest priority, ensuring rapid turnaround times and a responsive experience for the largest segment of the user base. 

UT2D was designed as a pair of autonomous operating systems that communicated to share resources like extended core memory, mass storage, and system libraries. The 6600 was primarily dedicated to heavy batch processing jobs, which demanded high computational power for scientific simulations. The 6400 was assigned to handle the interactive jobs and student programs that arrived through the timesharing system. UT2D was written primarily in Compass, the assembly language for the CDC 6000 series, and it remained a critical part of the university’s infrastructure into the mid eighties. The resulting infrastructure proved remarkably resilient, maintaining its utility for over two decades and continuing to support core accounting and billing functions into the nineties on the later CDC Cyber. The Dual Dinosaur era represents a unique period in high-performance computing where a small, dedicated staff could successfully challenge vendor dominance through local innovation. 

By the early seventies, the need for interactive time-sharing led to the development of the Taurus system, in parallel with UT2D. It allowed researchers at remote locations, such as the Marine Science Institute in Galveston, to program and access data files through teletype terminals, spreading access across the UT system. The Taurus design philosophy was captured in its full name, the Texas Anthropocentric Ubiquitous Responsive User System, signaling a deliberate focus on the human user experience during an era when computing was often machine-centric, and framing it as a philosophical counterpart to the student-focused scheduling of UT2D. Taurus supported a large variety of text-based terminals and more advanced graphic terminal devices to also help with extending the Center's reach directly to users across campus and the state-wide system. 

The value of the 6600 lay not just in its power but its location within an open academic setting. While supercomputers at the national labs were distant and oversubscribed, dedicated to high-priority government research projects, the university setting provided a more open and accessible resource for foundational research. The Computation Center intentionally cultivated outreach projects alongside the growth of its own institutional skills, with the result that a local technical culture developed from 1966 onward, creating a unique lineage of software. The scientific and operating system programming skills, specialized libraries, and expert personnel brought into place during this period served as the foundation for much that was to follow.

\section{Southwest Network And ARPANET}

In 1969 the Southwest Region Educational Computer Network was created around the Computation Center’s infrastructure, with funding from the National Science Foundation. By 1972 the network consisted of eleven four-year colleges and universities and three secondary schools across Texas. Computation Center staff and the computer science faculty provided the core resources for the network. The establishment of the network was a pioneering step in regional academic collaboration that cemented UT Austin’s status as the regional hub for computing and served as a precursor to modern distributed computing networks and the Internet. Of the fourteen institutions in the network, only UT Austin, Rice and Trinity universities had sizable computing facilities and well-established computing programs on campus. For the others, without adequate computing facilities and without sufficient funds to acquire them, remote computing services through the network offered a compromise solution. 

The primary objective was the modification and augmentation of the curricula of the network schools by adding computer science courses and by introducing computer methods into established programs. That required the implementation of the necessary computer terminals, communication links, central site computer hardware, and software systems and development of a library of computer programs to support curriculum development. All network schools except Rice University and Southwest Texas State University communicated with the UT2D system exclusively through teletypes. Rice University operated a CDC 200 User Terminal at 2000 baud in a remote job entry mode. Southwest Texas State University operated an IBM 1130, also in remote job entry mode. 

1969 also saw the beginnings of the national Arpanet, soon to grow into the Internet. While UT Austin was a significant early node, it was not one of the 1969 original four nodes, nor was it part of the network during its initial expansion in the early seventies. The UT in the original four 1969 Arpanet nodes referred to the University of Utah’s node four, which was installed in December 1969. The other three original 1969 nodes were UCLA, the Stanford Research Institute, and UC Santa Barbara. UT Austin does not appear on Arpanet logical maps from the 1969 to 1974 period. A management study report from January 1974 explicitly noted that major commercial and research centers in Texas were not yet represented on the network. 

The first IMP was installed at UT Austin in 1977. By the March 1977 Arpanet map, Texas is clearly visible as a node, and the Arpanet Directory from 1978 lists the University of Texas at Austin, specifically the Linguistics Research Center and Computation Center, as a fully operational node. In the transition to the modern Internet in the eighties, UT Austin was assigned significant network resources. The Class A network 39.0.0.0 was assigned to the UT Austin Balcones Research Center in later IP registries.

The CDC machines with their dual cathode ray tube display consoles were the roaring heart of campus technical life up through the mid seventies. Students submitted their programs on decks of punched cards at a number of remote job entry locations scattered across campus. An RJE accepted job decks and printed out and returned job results as they became available. A major RJE was located in Pearce Hall, the old law building, which was replaced by the new graduate school of business in the mid seventies. Another RJE was located in Taylor Hall, which eventually became the home of the computer science department during the eighties, after migrating over from Painter Hall. The reliability of the 6600 during this period was often a point of pride. Despite modern myths that the machines failed every few hours, researchers at the time recalled the 6600 running for months without unscheduled maintenance. This reliability allowed for marathon computation sessions, such as the calculation of Pi to 500,000 decimal places over a sixty-hour Thanksgiving weekend, a task that required periodic checkpoint dumps to save intermediate results, as the standard job control cards only allowed for about nine hours of requested job time.
